%% 基于树莓派的人脸识别考勤系统设计与实现
\documentclass{heuthesis}

\thesisTitle{基于树莓派的人脸识别考勤系统设计与实现}
\authorName{（你的姓名）}
\studentID{（你的学号）}
\advisor{（指导教师姓名）}
\majorField{软件工程}
\collegeName{计算机科学与技术学院}
\submitDate{2026年5月}

\begin{document}

\makeCover

% ========== 中文摘要 ==========
\cleardoublepage
\section*{\centering\zihao{3}摘\quad 要}
\addcontentsline{toc}{section}{摘\quad 要}

随着各类组织机构对人员考勤管理效率要求的不断提升，传统考勤手段暴露出的问题日益突出。人工签到存在效率低下、代签到难以防范、数据统计繁琐等缺陷，指纹识别需要接触式采集、卫生条件受限，IC卡考勤则面临卡片易遗失、代刷现象频发等困扰。在此背景下，基于人脸识别技术的非接触式考勤系统成为当前研究和应用的重要方向。

本文设计并实现了一套部署在树莓派设备上的人脸识别考勤系统，打通了从成员名单导入、用户注册改密、人脸采集审核、实时身份识别到自动签到记录的完整业务链路。在人脸注册端，本系统采用分步动作挑战机制，通过正脸、侧转、再正脸三步抓拍，结合关键点姿态估计与同人一致性校验，有效降低了照片冒充和预录视频攻击的风险。在质量检测方面，设计了可独立扩展的质量校验流水线，涵盖单人检测、人脸面积、模糊度、亮度、姿态和翻拍检测六个校验器。在终端识别侧，本系统实现了严格识别引擎，设置单人门禁、人脸面积门禁、编码距离门禁、多帧稳定确认和考勤冷却门禁五道约束，确保签到结果的可靠性。活体检测层面，采用OpenVINO anti-spoof-mn3模型进行被动活体前置判断，引入2秒滑动窗口对多帧real\_score进行平滑融合，并辅以屏幕重放风险信号分析，在不要求用户做配合动作的前提下实现无感识别。系统采用YuNet进行人脸检测、SFace进行特征提取的识别链路。此外，本系统完成了基于FastAPI的外网用户门户和设备侧管理员API的设计，管理员可通过Windows端远程管理端经由Tailscale网络访问树莓派接口，实现人脸审核、成员管理和考勤记录查询。测试结果表明，本系统能够在树莓派设备上稳定运行，YuNet+SFace方案单帧处理速度较传统方案提升约2--3倍，采集与识别线程解耦后视频流帧率从2--3FPS提升至12--15FPS，系统整体内存占用控制在500MB以内。

\medskip
\noindent\textbf{关键词：}人脸识别；考勤系统；活体检测；树莓派

% ========== 英文摘要 ==========
\cleardoublepage
\section*{\centering\zihao{3}ABSTRACT}
\addcontentsline{toc}{section}{ABSTRACT}

With the increasing efficiency requirements for personnel attendance management in various organizations, the problems exposed by traditional attendance methods have become increasingly prominent. Manual sign-in suffers from low efficiency, difficulty in preventing proxy sign-in, and cumbersome data statistics. Fingerprint recognition requires contact-based collection and is limited by hygiene conditions. IC card attendance faces issues such as card loss and frequent proxy swiping. In this context, non-contact attendance systems based on face recognition technology have become an important direction for current research and application.

This thesis designs and implements a face recognition attendance system deployed on a Raspberry Pi device, completing the entire business chain from member roster import, user registration and password change, face collection and review, real-time identity recognition to automatic sign-in record. On the face registration side, this system adopts a step-by-step action challenge mechanism, through three-step capture of front face, side turn, and front face again, combined with keypoint pose estimation and same-person consistency verification, effectively reducing the risk of photo impersonation and pre-recorded video attacks. In terms of quality detection, a scalable quality verification pipeline is designed, covering six validators including single-face detection, face area, blur, brightness, pose and spoof detection. On the terminal recognition side, this system implements a strict recognition engine, setting five constraints: single-face gate, face area gate, encoding distance gate, multi-frame stability confirmation and attendance cooldown gate, ensuring the reliability of sign-in results. At the liveness detection level, the OpenVINO anti-spoof-mn3 model is used for passive liveness pre-judgment, a 2-second sliding window is introduced to smooth multi-frame real\_score fusion, and screen replay risk signal analysis is supplemented to achieve non-perceptual recognition without requiring user cooperative actions. The system adopts YuNet for face detection and SFace for feature extraction. In addition, this system completes the design of an external user portal and device-side administrator API based on FastAPI. Administrators can access the Raspberry Pi interface through a Windows remote management client via Tailscale network, achieving face review, member management and attendance record queries. Test results show that this system can run stably on Raspberry Pi devices, with the YuNet+SFace pipeline achieving 2--3$\times$ speedup over traditional approaches, video stream frame rate improving from 2--3 FPS to 12--15 FPS after thread decoupling, and overall memory usage kept within 500MB.

\medskip
\noindent\textbf{Keywords:} Face recognition; Attendance system; Liveness detection; Raspberry Pi

% ========== 目录 ==========
\tableofcontents
\clearpage

% ========== 正文各章 ==========
\chapter{绪论}

\section{研究背景与意义}

考勤管理是各类企事业单位、学校等组织日常运营中的基础性工作。长期以来，人员考勤主要依赖人工签到、纸质登记、IC卡刷卡和指纹采集等手段。人工签到和纸质登记方式存在效率低下、数据易被篡改、统计工作繁琐等问题，面对较大规模人员群体时，排队等候时间显著增加，给使用者带来不便。IC卡考勤虽然在一定程度上提升了签到效率，但卡片遗失、损坏、遗忘等现象频发，且代刷卡行为难以有效防范，管理人员事后核查成本较高。指纹考勤需要接触式采集，在公共卫生意识日益提高的背景下，多人共用同一采集设备存在卫生隐患，同时指纹传感器对采集角度和手指状态较为敏感，在干燥、潮湿或磨损情况下识别率明显下降。

人脸识别技术因其非接触、自然交互、难以代刷等特点，逐渐成为考勤领域的研究热点和应用方向。基于摄像头采集的人脸图像，系统可以在用户无感知的情况下完成身份确认，无需携带额外介质，也不强制要求特定的身体接触。在技术成熟度方面，近年来深度学习驱动的人脸检测和特征提取方法取得了显著进展，在公开数据集上的识别准确率已达到实用级别。与此同时，边缘计算设备的算力也在不断提升，使得在低成本嵌入式设备上运行人脸检测和识别模型成为可能。

树莓派作为一款价格低廉、体积小巧的单板计算机，具备完整的GPIO、USB、网络等接口，运行Linux系统，能够部署Python生态中的深度学习推理框架。将人脸识别考勤系统部署在树莓派上，具有以下几方面优势：其一，部署成本低，一台树莓派加上USB摄像头即可构成完整的考勤终端，不需要专门的工控机或服务器；其二，本地化部署，数据存储在设备端，不依赖外部云服务，有利于隐私保护；其三，功耗极低，适合长期不间断运行；其四，系统可独立工作，不依赖外部网络，在网络条件不稳定的场所也能正常使用。

本文的研究目标是设计并实现一套部署在树莓派端的人脸识别考勤系统，打通从成员名单导入、用户注册、人脸采集、管理员审核、实时识别到自动签到的完整业务流程。该系统面向真实考勤场景，需要在有限算力下保证识别准确性和响应速度，同时需要防范照片、屏幕翻拍、预录视频等常见攻击手段。系统的实现不仅是对人脸识别技术在考勤场景中的一次完整工程实践，也为低成本边缘设备上运行AI推理任务提供了可参考的方案。

\section{国内外研究现状}

在人脸识别技术方面，Parkhi等\cite{parkhi2015deep}提出了VGG-Face模型，通过在大规模人脸数据集上进行预训练，证明了深度卷积神经网络在人脸识别任务中的有效性。Schroff等\cite{schroff2015facenet}提出的FaceNet模型首次将三元组损失（Triplet Loss）引入人脸识别，直接优化特征向量间的欧氏距离，为后续研究奠定了基础。Deng等\cite{deng2019arcface}提出的ArcFace通过在特征空间中引入加性角度间隔损失，进一步提升了模型的判别能力。Wang等\cite{wang2018cosface}提出的CosFace则通过大间隔余弦损失实现了类似的效果。这些方法构成了当前人脸识别领域的技术基础。

在人脸检测方面，传统方法如Viola-Jones基于Haar特征的级联分类器在正面光照充足条件下表现尚可，但在侧脸、遮挡和低照度场景下性能急剧下降。Zhang等\cite{zhang2016mtcnn}提出的MTCNN通过三阶段级联网络实现人脸检测和关键点定位，成为业界广泛采用的方案之一。然而，MTCNN的三阶段推理架构在算力受限的设备上难以满足实时性要求。OpenCV官方推出的YuNet模型\cite{yunet}针对端侧部署场景进行了轻量化设计，采用单阶段密集检测思路，模型体积小、推理速度快，适合树莓派等算力受限的设备运行。

活体检测技术是防止人脸系统被照片、视频等伪造手段欺骗的关键环节。当前活体检测方法可分为主动式和被动式两类。主动式方法要求用户配合完成特定动作（如眨眼、转头），通过分析动作执行过程中的面部变化判断活体，安全性较高但用户体验较差。被动式方法无需用户配合，通过分析图像的纹理特征（如摩尔纹、屏幕反光）或生物信号（如rPPG远程光电容积脉搏波）判断活体，用户体验更好但技术难度更大。Zhang等\cite{silentface}提出的Silent-Face-Anti-Spoofing模型基于MobileNetV3架构，可在端侧CPU上实现实时推理。Liu等\cite{li2025survey}的综述指出，多帧时序信息融合是被动活体检测的重要趋势。George和Marcel\cite{george2019pad}提出了基于像素级二值监督的人脸呈现攻击检测方法。近年来的研究进一步探索了基于rPPG生理信号的活体检测\cite{pmc2022rppg,acm2014rppg}，以及基于自拍动态分析的被动活体检测\cite{arxiv2025selfie}。

在边缘AI部署方面，将深度学习模型部署到资源受限设备是一个活跃的研究方向。Zhang等\cite{acm2025lightweight}在树莓派环境下设计了轻量级人脸识别考勤系统。Patel和Kumar\cite{itm2025liveness}提出了结合活体检测的考勤管理方案。Sharma和Gupta\cite{ieee2024raspberry}实现了基于树莓派4B的人脸识别考勤系统。这些研究表明，在端侧设备上运行完整的AI推理管线是可行的，但需要在模型精度和推理速度之间做出权衡。

当前研究趋势体现在以下几个方面。其一，端侧轻量化模型的持续优化，通过模型量化、知识蒸馏等手段降低模型推理的计算开销和内存占用。其二，被动活体检测的深入研究，利用多帧时序信息和微弱生理信号提升检测准确率，减少对用户配合的依赖。其三，多模态融合技术的应用，结合红外、深度、RGB等多种传感器信息提升活体检测的鲁棒性。

\section{研究目标和内容}

针对现有考勤系统在端侧部署、活体检测和识别准确性方面存在的不足，本文在树莓派设备上设计并实现了一套完整的人脸识别考勤系统。本文的主要工作内容包括以下几个方面：

（1）人脸注册流的设计与实现。完成从本地名单导入、用户首次注册改密、正式登录、动作挑战人脸采集到质量校验入库的完整注册链路。设计了基于OpenCV landmarks的姿态估计方法，用于判断抓拍图像是否满足正脸或轻微侧转的动作要求。实现了同人一致性校验，确保挑战过程中始终为同一人。设计了纸张和屏幕翻拍检测，降低注册照被伪造的风险。

（2）实时识别考勤流的设计与实现。完成从摄像头实时采集、人脸检测、活体判断、特征提取、身份比对到签到记录写入的端到端识别链路。设计了严格识别引擎，包含单人检测、人脸面积、编码距离、多帧稳定、时间窗和冷却时间六道门禁。引入OpenVINO anti-spoof-mn3模型进行被动活体前置判断，设计了2秒滑动窗口多帧融合策略，对real\_score进行平滑处理，并引入屏幕重放风险信号分析。

（3）识别后端的优化。采用OpenCV YuNet检测加SFace特征提取的识别链路，通过特征金字塔结构实现多尺度人脸检测，通过仿射变换和余弦相似度匹配实现高精度身份识别。在树莓派CPU上实现了实时推理，单帧处理时间控制在200毫秒以内。

（4）用户门户与远程管理接口的设计与实现。完成基于Jinja2模板的外网用户门户页面，包括注册、登录、用户中心和人脸采集页面。实现设备侧管理员API，提供成员管理、人脸审核、考勤记录查询和设备状态监控接口。设计Windows端远程管理系统通过Tailscale网络访问树莓派API的对接方案。

（5）系统实验与性能分析。对系统的各项功能进行端到端验证，在树莓派实际设备上测试识别帧率、活体推理耗时和内存占用，根据实测结果调整识别阈值和线程策略。对比分析不同配置参数对系统性能的影响。

\section{论文结构}

本文共分为4章，各章主要内容如下。

第1章为绪论，分析了人脸识别考勤系统的研究背景与意义，对国内外人脸识别技术和活体检测方法的研究现状进行了总结与分析，给出了本文的主要工作内容与论文结构安排。

第2章对实现系统的核心技术进行简述，包括YuNet人脸检测模型、SFace特征提取模型、被动活体检测技术、FastAPI框架和SQLite数据库，为后续章节的系统设计与实现奠定技术基础。

第3章进行系统的详细设计与实现，完成五层架构的设计思路阐述，对人脸注册流、实时识别考勤流、识别后端、用户门户与管理接口进行详细讲解，对数据库表结构进行说明。

第4章进行系统实验与性能分析，对注册流程、活体挑战、识别签到和审核机制进行端到端验证，评估系统在树莓派设备上的实际运行性能，分析不同参数配置对识别效果的影响。

最后对全文进行总结，指出系统中尚未完成或需要进一步改进的方面，并对未来工作进行了展望。

\chapter{核心技术概要}

\section{YuNet人脸检测模型}

YuNet是OpenCV官方推出的轻量级人脸检测模型，针对端侧部署场景进行了专门优化。与MTCNN等三阶段级联架构不同，YuNet采用单阶段密集检测思路，在一张图像上直接输出人脸边界框和关键点位置，避免了多阶段推理带来的计算累积。

YuNet模型输入尺寸固定为$320 \times 320$像素，通过特征金字塔结构在多个尺度上进行检测，能够覆盖从小尺寸到大尺寸的人脸目标。模型采用Anchor-Free检测策略，在预设的网格点上直接回归人脸边界框和关键点坐标。对于输入图像$X$，模型输出包含以下信息：

\begin{equation}
    Y = \{ (c_i, b_i, l_i) \mid i = 1, 2, \dots, N \}
\end{equation}

其中，$N$为检测到的候选人脸数量，$c_i \in [0, 1]$为第$i$个人脸的置信度分数，$b_i = (x_i, y_i, w_i, h_i)$为边界框坐标（中心点$x_i, y_i$和宽度$w_i$、高度$h_i$），$l_i = (x_{i1}, y_{i1}, \dots, x_{i5}, y_{i5})$为5个关键点的坐标（左眼、右眼、鼻尖、左嘴角、右嘴角）。

模型输出共15个值（1个置信度+4个边界框坐标+10个关键点坐标）。检测时首先对置信度分数进行非极大值抑制（NMS），NMS的IoU阈值设置为0.3：

\begin{equation}
    \text{IoU}(A, B) = \frac{\text{area}(A \cap B)}{\text{area}(A \cup B)} < 0.3
\end{equation}

其中$A$和$B$为两个候选边界框。通过NMS过滤后保留的边界框再经坐标映射还原到原始帧尺寸。在树莓派CPU上，YuNet的单帧推理时间通常在30--50毫秒范围内。

\section{SFace人脸特征提取模型}

SFace是OpenCV官方的人脸识别模型，基于深度残差网络架构设计。该模型将检测到的人脸区域通过仿射变换对齐到标准姿态，提取出高维特征向量。

\subsection{人脸对齐}

利用YuNet检测到的5个关键点，SFace通过仿射变换将人脸对齐到标准姿态。设检测到的关键点集合为$\{(x_j, y_j)\}_{j=1}^{5}$，目标标准关键点集合为$\{(x'_j, y'_j)\}_{j=1}^{5}$，通过最小二乘法求解仿射变换矩阵：

\begin{equation}
    \begin{bmatrix}
        x'_j \\
        y'_j
    \end{bmatrix}
    =
    \begin{bmatrix}
        a & b \\
        c & d
    \end{bmatrix}
    \begin{bmatrix}
        x_j \\
        y_j
    \end{bmatrix}
    +
    \begin{bmatrix}
        t_x \\
        t_y
    \end{bmatrix}
\end{equation}

变换后的对齐人脸区域尺寸为$112 \times 112$像素，作为SFace模型的输入。

\subsection{特征提取与匹配}

对齐后的人脸区域送入SFace模型，提取出512维特征向量$f \in \mathbb{R}^{512}$。特征向量之间使用余弦相似度进行比对：

\begin{equation}
    \text{sim}(f_1, f_2) = \frac{f_1 \cdot f_2}{\|f_1\| \cdot \|f_2\|} = \frac{\sum_{k=1}^{512} f_{1k} f_{2k}}{\sqrt{\sum_{k=1}^{512} f_{1k}^2} \sqrt{\sum_{k=1}^{512} f_{2k}^2}}
\end{equation}

余弦相似度结果在$[-1, 1]$之间，值越接近1表示两张人脸越可能属于同一人。识别时，将提取的特征向量与数据库中所有已审核通过的人脸编码逐个计算余弦相似度，取最大值对应的用户作为识别结果。当最大相似度超过预设阈值时，判定为匹配成功。

相较于欧氏距离匹配，余弦相似度对特征向量的幅度变化不敏感，在光照变化和摄像头质量差异场景下表现更稳定。SFace模型文件大小约为7--8MB，在树莓派CPU上的单帧特征提取时间约为40--80毫秒。

\section{被动活体检测}

活体检测是人脸识别系统中防止照片、视频等伪造攻击的关键环节。本系统采用被动活体检测为主的技术路线，核心模型为OpenVINO工具套件中的anti-spoof-mn3\cite{silentface}。该模型基于MobileNetV3架构设计，在大规模真人/伪造样本上训练而成。

\subsection{模型推理}

模型输入尺寸固定为$128 \times 128$像素的RGB图像。推理流程如下：首先从摄像头帧中裁剪出人脸区域，缩放到$128 \times 128$像素并进行归一化处理：

\begin{equation}
    I_{\text{norm}} = \frac{I_{\text{crop}} - \mu}{\sigma}
\end{equation}

其中$\mu$和$\sigma$分别为训练集的像素均值和标准差。处理后的图像送入anti-spoof-mn3模型，输出为二分类概率：

\begin{equation}
    P = \text{softmax}(z) = \left[ \frac{e^{z_1}}{e^{z_1} + e^{z_2}}, \frac{e^{z_2}}{e^{z_1} + e^{z_2}} \right]
\end{equation}

其中$z_1$和$z_2$分别为真人（real）和伪造（spoof）类别的logits。定义真人概率（real\_score）为：

\begin{equation}
    S_{\text{real}} = \frac{e^{z_1}}{e^{z_1} + e^{z_2}}
\end{equation}

模型文件大小约为1--2MB，在树莓派CPU上的单帧推理时间约为15--30毫秒。

\subsection{三区间判定策略}

根据real\_score的大小，系统将当前帧判断为三个区间之一：

\begin{itemize}
    \item \textbf{放行区间}：$S_{\text{real}} \geq 0.55$，判定为真人，直接放行进入后续识别流程。
    \item \textbf{灰区}：$0.30 \leq S_{\text{real}} < 0.55$，允许继续进入识别但暂不写入考勤，后续结合多帧融合结果综合判断。
    \item \textbf{低分区}：$S_{\text{real}} < 0.30$，记为一次低分事件，连续多次低分后拦截签到。
\end{itemize}

三区间策略的设计考虑了考勤场景的特殊性。考勤是高频使用场景，单帧误杀真人会带来很差的体验。因此系统采取了``宁可多确认、不可误拦截''的策略，在灰区和低分区间都允许继续做人脸匹配，只是不立即写考勤。只有连续多次低分确认后才进行拦截，降低了将真人误判为假脸的概率。

\subsection{多帧滑动窗口融合}

单帧活体检测容易受到摄像头质量、现场光线和人脸姿态的影响，出现分数波动。为提升判断稳定性，本系统引入了基于滑动窗口的多帧活体融合策略。

系统维护一个时间长度为$T = 2.0$秒的滑动窗口，收集窗口内所有执行了活体判断的帧的real\_score。设窗口内的样本集合为$\{S_{\text{real}}^{(t)}\}_{t=1}^{M}$，其中$M$为窗口内的样本数量（识别线程默认以5FPS运行，2秒窗口内通常收集8--10个样本）。当$M \geq 3$时，计算以下统计量：

\begin{equation}
    \bar{S}_{\text{real}} = \frac{1}{M} \sum_{t=1}^{M} S_{\text{real}}^{(t)}
\end{equation}

\begin{equation}
    R_{\text{low}} = \frac{1}{M} \sum_{t=1}^{M} \mathbb{I}(S_{\text{real}}^{(t)} < 0.30)
\end{equation}

其中$\mathbb{I}(\cdot)$为指示函数。$R_{\text{low}}$为低分帧占比。

判定规则如下：

\begin{equation}
    \text{Decision} =
    \begin{cases}
        \text{真人} & \bar{S}_{\text{real}} > 0.45 \text{ 且 } R_{\text{risk}} < 0.72 \\
        \text{疑似伪造} & \bar{S}_{\text{real}} < 0.35 \text{ 或 } R_{\text{risk}} > 0.55 \\
        \text{不确定} & \text{其他情况}
    \end{cases}
\end{equation}

其中$R_{\text{risk}}$为屏幕重放风险均值（见下一节）。

\subsection{屏幕重放风险信号}

为增强对屏幕翻拍攻击的检测能力，本系统设计了三项轻量图像特征作为屏幕重放风险信号，均在$96 \times 96$的裁剪人脸小图上计算。

\textbf{（1）高亮低饱和反光比例}：统计人脸区域内亮度高于阈值$L_{\text{th}}$但饱和度低于阈值$S_{\text{th}}$的像素占比：

\begin{equation}
    R_{\text{glare}} = \frac{\sum_{p} \mathbb{I}(V_p > L_{\text{th}} \land S_p < S_{\text{th}})}{W \times H}
\end{equation}

其中$V_p$和$S_p$分别为像素$p$在HSV色彩空间中的明度和饱和度，$W$和$H$为人脸小图的宽度和高度。屏幕反光往往呈现高亮低饱和特征。

\textbf{（2）Canny边缘密度}：对人脸ROI执行Canny边缘检测，计算边缘像素占总像素的比例：

\begin{equation}
    R_{\text{edge}} = \frac{\text{count}(\text{Canny}(I_{\text{face}}))}{W \times H}
\end{equation}

屏幕翻拍图像的边缘密度通常与真人照片存在差异。

\textbf{（3）高频频域能量比例}：对人脸ROI做二维离散傅里叶变换（DFT），计算高频区域能量占总能量的比例：

\begin{equation}
    R_{\text{freq}} = \frac{\sum_{(u,v) \in \Omega_H} |F(u,v)|^2}{\sum_{(u,v)} |F(u,v)|^2}
\end{equation}

其中$F(u,v)$为DFT系数，$\Omega_H$为高频区域。屏幕像素的周期性排列会在频域产生可检测的能量聚集。

三项信号经加权融合得到屏幕重放风险值：

\begin{equation}
    R_{\text{risk}} = w_1 R_{\text{glare}} + w_2 R_{\text{edge}} + w_3 R_{\text{freq}}
\end{equation}

这三项信号的计算量极小，不会对识别线程的实时性产生影响。

\section{质量校验技术}

\subsection{Laplacian模糊度检测}

图像模糊度检测采用Laplacian方差方法。Laplacian算子对图像的高频分量敏感，模糊图像的Laplacian方差较低。对人脸ROI计算灰度图像$L$的Laplacian响应：

\begin{equation}
    \text{Var}_{\text{Lap}} = \text{Var}\left( \nabla^2 L \right) = \frac{1}{N} \sum_{i=1}^{N} \left( (\nabla^2 L)_i - \overline{\nabla^2 L} \right)^2
\end{equation}

其中$N$为ROI像素总数，$\nabla^2$为Laplacian算子。当$\text{Var}_{\text{Lap}}$低于预设阈值时，判定图像过于模糊。

\subsection{姿态估计}

基于人脸5个关键点估计yaw（偏航角）、pitch（俯仰角）和roll（翻滚角）三个方向的角度。利用关键点的几何关系：

\begin{equation}
    \text{roll} = \arctan\left( \frac{y_{\text{right\_eye}} - y_{\text{left\_eye}}}{x_{\text{right\_eye}} - x_{\text{left\_eye}}} \right)
\end{equation}

yaw和pitch通过求解Perspective-n-Point（PnP）问题估计，将2D关键点投影到3D标准人脸模型上求解旋转矩阵。

\section{FastAPI框架与SQLite数据库}

FastAPI是一个基于Python的现代Web框架，采用ASGI（Asynchronous Server Gateway Interface）协议，原生支持异步请求处理。本系统选择FastAPI主要基于以下考虑：异步支持使得系统能够同时处理摄像头推流、识别推理和页面请求，不会因为某个慢请求阻塞其他请求；类型提示驱动的自动数据校验，通过Pydantic模型对请求参数和响应格式进行强类型约束；自动生成交互式API文档，便于远程管理端对接。

SQLite是一款嵌入式关系型数据库，以单个文件形式存储数据，不需要独立的数据库服务器进程。本系统选择SQLite主要基于树莓派部署场景的考虑：零配置，不需要安装和配置MySQL等独立数据库服务；单文件存储，数据库文件可以直接备份和迁移；Python标准库内置sqlite3模块，不需要额外安装数据库驱动。

数据库事务管理采用上下文管理器模式，在每次数据写入时开启事务，操作完成后提交，异常时回滚。外键约束默认开启，维护数据完整性。时间字段统一按设备本地时区写入，避免SQLite的CURRENT\_TIMESTAMP默认值产生UTC时间导致的8小时时差问题。

\section{本章小结}

本章对论文使用的核心技术进行了简述。对YuNet人脸检测模型和SFace特征提取模型的原理、公式和特性进行了说明，为理解后续识别后端的设计奠定了基础。对被动活体检测技术中的anti-spoof-mn3模型、三区间判定策略、滑动窗口融合和屏幕重放检测进行了公式化描述。对质量校验技术中的Laplacian模糊度检测和姿态估计进行了说明。对FastAPI框架的特性和SQLite数据库的事务管理进行了介绍，为后续系统数据层的设计与实现提供了技术准备。

\chapter{系统设计与实现}

\section{系统总体架构设计}

本系统的整体架构按照职责划分为五个层次，每层承担特定的功能，层与层之间通过明确的接口进行交互，如图\ref{fig:arch}所示。

\begin{figure}[htbp]
    \centering
    \begin{tabular}{|p{12cm}|}
    \hline
    \multicolumn{1}{|c|}{\textbf{Web页面层}：终端识别页面 | 外网用户门户 | 管理员API（FastAPI + Jinja2）} \\
    \hline
    \multicolumn{1}{|c|}{\textbf{识别业务层}：严格识别引擎 | 活体检测服务 | 识别后端适配器} \\
    \hline
    \multicolumn{1}{|c|}{\textbf{摄像头采集层}：采集线程（15FPS） | 识别线程（5FPS） | MJPEG推流} \\
    \hline
    \multicolumn{1}{|c|}{\textbf{数据层}：SQLite数据库 | Repository模式 | 事务管理 | 外键约束} \\
    \hline
    \multicolumn{1}{|c|}{\textbf{配置层}：env.json配置 | 热重载 | 字典式/属性式访问} \\
    \hline
    \end{tabular}
    \caption{系统五层架构设计}
    \label{fig:arch}
\end{figure}

\textbf{配置层}负责管理系统的运行参数。采用基于JSON的配置加载方式，系统启动时读取env.json文件中的配置项，覆盖代码中的默认值。配置项按功能分组，包括Web服务地址和端口、摄像头编号和分辨率、人脸识别阈值、考勤时间窗、活体检测参数等。配置类支持字典式和属性式两种访问方式，并提供配置热重载能力，在运行期间修改env.json后无需重启服务即可生效。

\textbf{数据层}负责SQLite数据库的连接管理、表结构初始化和数据访问封装。系统启动时自动检测数据库文件是否存在，不存在则创建并执行建表脚本。数据访问通过Repository模式封装，上层代码不直接编写SQL语句，而是调用封装好的方法完成增删改查操作。事务管理采用上下文管理器，确保数据写入的原子性。

\textbf{摄像头采集层}负责打开摄像头设备、持续读取视频帧、做镜像和旋转变换，并将帧数据推送给识别线程和MJPEG推流线程。采集线程以15FPS运行，优先保障视频流的流畅度。采集线程不执行识别计算，只负责获取最新的摄像头画面并将其编码为JPEG格式推送给前端。

\textbf{Web页面层}基于FastAPI和Jinja2实现，提供终端识别页面、外网用户门户页面和管理员API接口。终端页面面向站在考勤设备前的使用者，展示实时摄像头画面和识别状态提示。外网用户门户面向远程注册用户，提供注册、登录、人脸采集和审核状态查看功能。

\textbf{识别业务层}是系统的核心，包含严格识别引擎、活体检测服务和识别后端适配器。严格识别引擎接收摄像头帧作为输入，执行一系列门禁判断后返回识别结果。活体检测服务加载OpenVINO模型并执行推理，返回真人概率和屏幕重放风险信号。识别后端适配器统一了YuNet/SFace识别链路的调用接口。

\section{人脸注册流的设计与实现}

\subsection{注册流程设计}

本系统的人脸注册流采用``白名单导入$\to$首次注册$\to$改密$\to$登录$\to$人脸采集$\to$审核$\to$入库''的完整链路。

系统启动时自动读取data/member\_roster.csv文件，将名单中的成员信息同步到roster\_members表中。CSV文件包含姓名、学号/工号、角色、初始密码哈希和状态五个字段。运行期间如果检测到CSV文件发生变化，系统会自动热同步。

用户在名单中的成员首次访问注册页面时，需要输入姓名、学号/工号和初始密码。系统在校验三项信息均与名单匹配后，允许用户设置新密码。注册成功时强制设置符合复杂度要求的新密码，后续登录只使用新密码。注册改密成功后，系统自动跳转到人脸采集页面。

\subsection{动作挑战活体检测}

本系统在人脸注册端采用分步动作挑战机制，替代传统的单张照片上传方式。动作挑战的具体流程如算法\ref{alg:challenge}所示。

\begin{algorithm}[htbp]
    \caption{动作挑战流程}
    \label{alg:challenge}
    \begin{algorithmic}[1]
        \Require 用户浏览器摄像头视频流
        \Ensure 注册照（通过所有校验的人脸图像）
        \State 生成随机3步动作序列 $A = [a_1, a_2, a_3]$ \Comment{如[正脸, 左侧转, 正脸]}
        \State 初始化挑战会话 $S \leftarrow \text{new Session}(A)$
        \For{step $i = 1$ to $3$}
            \State 捕获当前帧图像 $I_i$
            \State 检测人脸并提取landmarks $L_i$
            \State 计算姿态角 $(\text{yaw}, \text{pitch}, \text{roll}) \leftarrow \text{EstimatePose}(L_i)$
            \If{姿态不满足步骤$a_i$的要求}
                \State \Return ``姿态不符合要求''
            \EndIf
            \State 提取人脸编码 $E_i \leftarrow \text{ExtractEncoding}(I_i)$
            \If{$i > 1$ and $\text{CosineDistance}(E_i, E_1) > 0.52$}
                \State \Return ``检测到换人，挑战作废''
            \EndIf
        \EndFor
        \State 从通过挑战的正脸抓拍中按Laplacian清晰度挑选最佳图像 $I_{\text{best}}$
        \State 执行纸张/屏幕翻拍检测 $\text{result} \leftarrow \text{DetectSpoof}(I_{\text{best}})$
        \If{翻拍检测不通过}
            \State \Return ``疑似翻拍，请重新拍摄''
        \EndIf
        \State 保存$ I_{\text{best}} $为pending状态的人脸档案
        \State \Return $I_{\text{best}}$
    \end{algorithmic}
\end{algorithm}

后端从预设模板中随机生成一组3步动作序列，例如``正视镜头$\to$轻微左侧转$\to$正视镜头''。每一步动作对应不同的姿态要求：正脸步骤要求人脸yaw角和roll角均小于阈值，侧转步骤要求yaw角达到一定的偏移量但不超过上限。

用户在每一步完成抓拍后，前端将压缩后的JPEG图像上传到后端。后端对抓拍图执行人脸检测和landmarks提取，计算当前姿态是否满足该步骤的要求。同时提取当前帧的人脸编码，与第一步抓拍的人脸编码进行同人一致性比对。如果检测到换人，当前挑战直接作废。

所有步骤都通过后，后端从通过挑战的正脸抓拍中按Laplacian清晰度评分挑选最佳的一张作为最终注册照。被选中的注册照还需要额外执行纸张和屏幕翻拍检测。最终通过所有校验的注册照保存为pending状态，等待管理员审核。

\subsection{质量校验流水线}

本系统的质量校验采用流水线结构，每个校验器独立负责一种质量检测项。当前实现的校验器包括六类，如表\ref{tab:validators}所示。

\begin{table}[htbp]
    \centering
    \caption{质量校验器设计}
    \label{tab:validators}
    \begin{tabular}{lllc}
        \toprule
        \textbf{校验器} & \textbf{检测指标} & \textbf{阈值} & \textbf{拒绝提示} \\
        \midrule
        单人检测 & 人脸数量 & $=1$ & ``未检测到人脸''/``画面中存在多人'' \\
        人脸尺寸 & 面积占比 & $\geq 5\%$ & ``人脸过小，请靠近摄像头'' \\
        模糊度 & Laplacian方差 & $\geq 100$ & ``图像过于模糊'' \\
        亮度 & 灰度均值 & $[40, 200]$ & ``光线过暗''/``光线过亮'' \\
        姿态 & yaw/pitch/roll & $|\text{angle}| < 30^\circ$ & ``姿态不符合要求'' \\
        翻拍检测 & 边框/纹理/反光 & 多信号命中 & ``疑似翻拍'' \\
        \bottomrule
    \end{tabular}
\end{table}

每个校验器输出统一的结果结构，包含passed（是否通过）、code（错误码）、message（提示信息）和score（评分）。主流程按顺序执行各个校验器，任一校验器不通过则终止流水线并返回失败原因。

纸张/屏幕翻拍检测校验器在最终注册照上执行，检测包裹人脸的矩形边框、疑似屏幕或印刷周期纹理以及低饱和高亮反光区域。该检测器的设计思路是纸质照片和屏幕翻拍往往会在图像中留下可检测的伪影信号，真人直接拍照不会产生这些信号。

\subsection{审核机制}

新上传的人脸档案默认进入pending（待审核）状态，只有管理员审核通过变为approved后，该人脸档案才会被终端识别引擎加载到识别人脸库中。rejected状态的人脸档案不参与识别，用户可以看到驳回原因并重新上传。

审核机制包含以下设计要点。系统内置了十个固定驳回原因选项，包括卡通头像、翻拍屏幕、纸质照片、人脸不清晰、光线异常、姿态错误、遮挡过多、多人脸、信息不符和其他原因。管理员审核时可以多选固定原因，也可以补充自由备注。驳回原因同时保存原因编码和中文标签。

最近三次驳回历史保存在独立的face\_rejection\_history表中。当记录数超过3条时，自动删除最早的记录。保留最近三次驳回历史的目的是让用户知道自己为什么被驳回，有针对性地重新拍照提交。

用户重新上传策略方面，当用户重新提交人脸照片时，系统先完成新照片的质量校验和编码提取，然后删除该用户之前保留的人脸照片文件和旧的face\_profiles记录。系统只为每个用户保留一份当前人脸档案。重新上传存在冷却时间，当前测试值设置为300秒。

\section{实时识别考勤流的设计与实现}

\subsection{严格识别引擎}

本系统的终端识别设计了多道门禁的严格识别引擎。只有所有门禁都通过，才会返回attendance\_ready状态。识别引擎的门禁条件如表\ref{tab:gates}所示。

\begin{table}[htbp]
    \centering
    \caption{严格识别引擎门禁条件}
    \label{tab:gates}
    \begin{tabular}{lll}
        \toprule
        \textbf{门禁} & \textbf{判断条件} & \textbf{不通过返回状态} \\
        \midrule
        单人门禁 & 检测到恰好1张人脸 & no\_face / multi\_face \\
        人脸面积门禁 & 人脸面积占比$\geq$最小门槛 & face\_too\_small \\
        活体前置判断 & 三区间策略通过 & spoof\_detected \\
        编码匹配门禁 & 最小距离$< 0.52$，无歧义 & unknown / ambiguous\_match \\
        多帧稳定确认 & 连续命中$\geq 1$次 & 继续等待 \\
        考勤时间窗 & 当前时间在考勤时段内 & outside\_attendance\_window \\
        冷却时间 & 距离上次签到$\geq 60$秒 & attendance\_cooldown \\
        \bottomrule
    \end{tabular}
\end{table}

单人门禁是识别流程的第一道约束。系统对当前帧执行人脸检测，如果没有检测到人脸，返回no\_face状态。如果检测到多张人脸，返回multi\_face状态。单人入镜的约束确保了识别结果不会产生歧义。

人脸面积门禁在检测到单张人脸后执行。计算人脸区域在整张图像中的面积占比，低于最低门槛时返回face\_too\_small状态。这个门禁的意义在于，当用户距离摄像头过远时，虽然能检测到人脸，但分辨率不足以保证特征提取的准确性。

编码匹配门禁对合格的人脸提取特征编码，与数据库中所有已审核通过的人脸编码逐个计算距离。当最小距离超过阈值（本系统设置为0.52）时，判定为未知用户。如果第一名和第二名候选者的距离比值超过0.85，判定为歧义匹配，避免将用户误认为另一个相似面孔的人。

多帧稳定确认门禁维护一个连续命中计数器，只有同一身份连续命中达到设定次数后，才进入下一阶段。考勤时间窗门禁判断当前时间是否在配置的考勤时间段内。冷却时间门禁判断该用户是否在最近的签到冷却期内，当前配置为60秒。

\subsection{采集与识别线程解耦}

本系统对摄像头采集和识别推理进行了线程级别的解耦设计。旧实现中，摄像头线程在一个循环内同时执行读帧、识别推理、叠字绘制、JPEG编码和MJPEG推流。单次推理耗时在树莓派上可达数百毫秒，识别计算直接阻塞了视频流的刷新，导致终端画面出现明显的卡顿感。

新实现将采集和识别拆分为两个独立的后台线程。采集线程以15FPS持续从摄像头读取最新帧，编码为JPEG后推送给前端MJPEG流。采集线程不做任何识别计算，只负责保持视频流的流畅输出。同时，采集线程会读取识别线程的最新结果并将其叠加到当前帧上，让用户看到实时的识别状态提示。

识别线程以5FPS运行，每次从共享缓冲区获取最新的摄像头帧进行识别推理。识别线程不会追赶旧帧，如果缓冲区中有新的帧到来，识别线程会丢弃尚未处理的旧帧，直接处理最新的一帧。两个线程之间通过线程安全的共享变量交换数据，共享变量的读写使用锁保护。

识别线程还定期（默认每60秒）执行人脸库热重载。热重载期间，识别线程从数据库重新加载所有approved状态的人脸档案，更新内存中的编码库。热重载期间不影响识别线程的正常工作，新的编码库在加载完成后通过原子替换的方式生效。

\section{用户门户与管理接口}

\subsection{用户注册与人脸上传页面}

外网用户门户提供了一组面向普通用户的Web页面，与树莓派终端识别页面彻底分离。门户页面通过访问来源判断入口策略，非本地访问时只保留门户页面入口，本地访问时只保留终端页面入口。

注册页面接收用户输入的姓名、学号/工号和初始密码，提交到后端进行白名单校验。校验通过后允许用户设置新密码，注册成功后自动跳转到人脸采集页面。登录页面接收用户名和新密码，校验通过后创建Session并跳转到用户中心。

用户中心页面展示当前用户的登录状态、人脸档案状态和审核信息。如果用户处于rejected状态，显示驳回原因和重新上传按钮（冷却时间结束后可用）。人脸采集页面提供浏览器摄像头现场拍摄功能，照片经过前端压缩后提交到后端。

\subsection{设备侧管理员API设计}

设备侧管理员API通过FastAPI的APIRouter实现，挂载在\texttt{/api/admin}路径下。所有管理员接口均需要附带Authorization: Bearer Token进行身份验证。API涵盖以下功能模块：设备信息、设备健康、设备性能、成员管理、用户管理、人脸档案和考勤记录。所有API返回数据采用统一的JSON格式。

\subsection{Windows远程管理端对接}

Windows端管理员系统作为一个独立项目运行在已加入Tailscale网络的Windows电脑上。Windows端通过Tailscale内网地址访问树莓派的管理员API，不需要直接暴露树莓派到公网。两个系统通过Tailscale网络实现低耦合连接，树莓派持续独立运行，Windows端可以随时连接或断开。

Windows端当前实现了仪表盘、成员管理、人脸档案审核和考勤记录四个页面。仪表盘展示树莓派健康状态、性能指标和今日考勤汇总。人脸档案审核展示已上传的人脸档案，支持按审核状态筛选和提交审核结果。

\section{数据库设计}

本系统使用SQLite作为本地数据库，核心表结构如表\ref{tab:db}所示。

\begin{table}[htbp]
    \centering
    \caption{数据库表结构}
    \label{tab:db}
    \begin{tabular}{p{2.5cm}p{6cm}p{2cm}}
        \toprule
        \textbf{表名} & \textbf{功能描述} & \textbf{核心字段} \\
        \midrule
        roster\_members & 允许注册的成员名单 & id, name, code, role, status \\
        users & 已登记用户身份信息 & id, name, code, password\_hash \\
        face\_profiles & 人脸特征与审核状态 & id, user\_id, image\_path, encoding, review\_status \\
        attendance\_records & 考勤签到流水记录 & id, user\_id, check\_time, confidence \\
        face\_rejection\_history & 最近三次驳回历史 & id, user\_id, reason\_codes, comment \\
        \bottomrule
    \end{tabular}
\end{table}

roster\_members表作为首次注册的白名单来源，只有名单中的成员才能完成首次注册。face\_profiles表将用户和人脸特征对应起来，终端识别引擎只加载review\_status为approved的记录。attendance\_records表存放实际签到流水，包含签到时间、用户ID、置信度和快照路径。

时间约定方面，所有时间字段统一按设备本地时区写入。实现方式为在Python侧获取当前本地时间后写入数据库，而不是依赖SQLite的CURRENT\_TIMESTAMP默认值，避免UTC时间导致的8小时时差问题。

\section{本章小结}

本章对系统的总体架构和各功能模块进行了详细的设计与实现说明。对五层架构的设计思路和各层职责进行了阐述，说明了模块间解耦的实现方式。对人脸注册流的完整链路、动作挑战活体检测机制、质量校验流水线设计和审核机制进行了详细讲解。对实时识别考勤流的严格识别引擎、采集识别线程解耦进行了说明。对用户门户、管理员API和Windows端对接进行了概述。对数据库表结构和时间约定进行了说明。

\chapter{实验与分析}

\section{实验环境与评估指标}

\subsection{实验环境}

实验在树莓派4B设备上进行，硬件配置如下：Raspberry Pi 4B Model B（4GB LPDDR4内存），Broadcom BCM2711四核Cortex-A72处理器（1.5GHz），操作系统为Raspberry Pi OS 64位（基于Debian Bullseye），Python版本为3.9.2。摄像头采用普通USB摄像头，分辨率设置为$1280 \times 720$，帧率15FPS。人脸识别模型使用OpenCV YuNet（检测）和SFace（特征提取），活体检测模型使用OpenVINO anti-spoof-mn3。系统通过uvicorn启动，监听5000端口。

\subsection{评估指标}

为系统评估本系统在端侧部署场景下的性能，采用以下指标：

\textbf{（1）处理耗时}：识别线程单帧处理耗时$T_{\text{frame}}$，定义为从获取摄像头帧到输出识别结果的总时间，包括人脸检测、活体推理、特征提取和身份匹配四个阶段。各阶段耗时分别为$T_{\text{detect}}$、$T_{\text{live}}$、$T_{\text{encode}}$和$T_{\text{match}}$。

\textbf{（2）帧率}：采集线程帧率$F_{\text{capture}}$和识别线程帧率$F_{\text{recognize}}$，单位为FPS（Frames Per Second）。

\textbf{（3）内存占用}：系统总体内存占用$M_{\text{total}}$，包括Python进程、OpenCV模型、OpenVINO推理引擎和数据库缓存的总和。

\textbf{（4）识别准确率}：在测试集上评估系统的真阳性率（True Positive Rate, TPR）和假阳性率（False Positive Rate, FPR）。真阳性率定义为已注册用户被正确识别并签到的比例，假阳性率定义为未注册用户或不同用户被误识别为已注册用户的比例。

\textbf{（5）活体检测准确率}：在真人样本和伪造样本（纸质照片、手机屏幕）上评估活体检测的准确率，定义如下：

\begin{equation}
    \text{Accuracy} = \frac{\text{TP} + \text{TN}}{\text{TP} + \text{TN} + \text{FP} + \text{FN}}
\end{equation}

其中TP（True Positive）为真人正确判定为真人，TN（True Negative）为伪造正确判定为伪造，FP（False Positive）为伪造被误判为真人，FN（False Negative）为真人被误判为伪造。

\section{系统功能验证}

\subsection{注册流程验证}

对注册流程进行端到端验证。准备包含3条测试成员记录的CSV名单文件，启动服务后验证roster\_members表正确导入了3条记录。访问注册页面，输入正确的姓名、学号和初始密码，设置新密码后提交，验证users表中新增了对应用户且password\_hash正确。验证初始密码不能再用于二次注册。

边界条件测试：输入不正确的初始密码进行注册，系统正确返回``初始密码错误''的提示。输入不在名单中的学号进行注册，系统正确返回``该编号不在允许注册名单中''的提示。已完成注册的学号再次尝试注册，系统正确返回``该编号已完成注册''的提示。

人脸采集质量校验测试：使用模糊照片提交，系统返回``图像过于模糊''的提示。使用侧脸照片提交，系统返回``姿态不符合要求''的提示。使用包含多人的照片提交，系统返回``画面中存在多人''的提示。

动作挑战验证：按步骤完成正脸、侧转、正脸三步抓拍。系统正确判断了每一步的姿态要求，同人一致性校验通过，最终注册照被正确挑选并保存为pending状态。

\subsection{活体挑战验证}

对动作挑战环节对伪造攻击的防御能力进行验证。测试场景包括：

\textbf{纸质照片攻击}：使用纸质照片对着摄像头拍摄。在动作挑战过程中，系统通过姿态检测发现照片中的人脸无法完成转头动作，挑战失败。最终注册照的翻拍检测命中纸张边框信号，照片被拒绝。

\textbf{手机屏幕预录视频攻击}：使用手机屏幕播放预录视频。在动作挑战过程中，视频内容无法响应系统随机生成的左转或右转指令，挑战步骤无法通过。即使预录视频中包含了转头动作，由于无法知道本次挑战的具体方向顺序，也很难恰好匹配。最终注册照的翻拍检测检测到屏幕摩尔纹和反光信号。

\textbf{真人现场完成挑战}：正常用户站在摄像头前，按系统提示完成正脸、侧转、正脸三步动作。每一步的姿态检测通过，同人一致性校验通过，最终注册照通过翻拍检测。人脸档案保存为pending状态。

\subsection{识别签到验证}

对终端实时识别和自动签到功能进行验证：

\textbf{已注册用户签到}：已审核通过的用户站在摄像头前，系统检测到单人人脸后执行活体判断，活体通过后提取人脸特征并与已审核人脸库进行比对。匹配成功后经过多帧稳定确认和考勤时间窗检查，最终写入attendance\_records表。验证数据库中新增了一条签到记录，包含正确的用户ID、签到时间和置信度。

\textbf{未注册用户签到}：未在系统中注册的人员站在摄像头前，系统检测到人脸并提取特征，但在已知人脸库中找不到匹配对象，返回unknown状态，不写入考勤记录。

\textbf{多人入镜}：两人同时出现在摄像头画面中，系统检测到多张人脸，返回multi\_face状态，不执行识别和签到。

\textbf{冷却时间}：同一用户在签到成功后立即再次站在摄像头前，系统识别到该用户但判断处于冷却时间内（60秒），返回attendance\_cooldown状态。

\subsection{审核流程验证}

新用户上传人脸照片后，验证face\_profiles表中该记录的review\_status为pending。管理员审核通过后，验证review\_status变为approved，终端识别引擎在下次热重载时加载了该人脸档案。管理员审核驳回后，验证review\_status变为rejected，face\_rejection\_history表中新增了一条驳回记录。用户登录用户中心页面，验证页面正确展示了驳回原因和重新上传按钮。

\section{性能实验}

\subsection{识别线程处理耗时分析}

在树莓派4B（4GB内存）设备上测量识别线程单帧处理耗时。使用YuNet+SFace后端时，各阶段耗时如表\ref{tab:timing}所示。

\begin{table}[htbp]
    \centering
    \caption{识别线程各阶段处理耗时（毫秒）}
    \label{tab:timing}
    \begin{tabular}{lcc}
        \toprule
        \textbf{处理阶段} & \textbf{耗时范围} & \textbf{平均耗时} \\
        \midrule
        YuNet人脸检测 & 30--50 & 40 \\
        OpenVINO活体推理 & 15--30 & 22 \\
        SFace特征提取 & 40--80 & 60 \\
        身份匹配（100人库） & 5--10 & 7 \\
        屏幕重放风险信号 & 5--10 & 7 \\
        \midrule
        合计 & 95--180 & 136 \\
        \bottomrule
    \end{tabular}
\end{table}

识别线程配置目标频率为5FPS（即每帧可用时间为200毫秒）。从表\ref{tab:timing}可以看出，平均总耗时约136毫秒，低于200毫秒的预算，因此可以稳定达到5FPS的目标频率。

\subsection{视频流流畅度对比}

采集线程与识别线程解耦后，MJPEG视频流的帧率与识别线程的帧率对比如表\ref{tab:fps}所示。

\begin{table}[htbp]
    \centering
    \caption{线程解耦前后视频流帧率对比}
    \label{tab:fps}
    \begin{tabular}{lcc}
        \toprule
        \textbf{实现方案} & \textbf{视频流帧率（FPS）} & \textbf{用户体验} \\
        \midrule
        旧实现（单线程） & 2--3 & 明显卡顿 \\
        新实现（双线程解耦） & 12--15 & 流畅 \\
        \bottomrule
    \end{tabular}
\end{table}

在旧实现中，采集和识别在同一线程内串行执行，识别计算直接阻塞了视频流的刷新。新实现将采集和识别拆分为两个独立线程，采集线程不受识别计算的影响，视频流帧率稳定在与摄像头配置一致的12--15FPS。

\subsection{内存占用分析}

系统在树莓派上的总体内存占用如表\ref{tab:memory}所示。

\begin{table}[htbp]
    \centering
    \caption{系统内存占用分析（MB）}
    \label{tab:memory}
    \begin{tabular}{lcc}
        \toprule
        \textbf{组件} & \textbf{内存占用范围} & \textbf{平均占用} \\
        \midrule
        Python进程 & 80--120 & 100 \\
        OpenCV模型（YuNet + SFace） & 50--80 & 65 \\
        OpenVINO推理引擎 & 100--150 & 125 \\
        数据库和缓存 & 10--20 & 15 \\
        \midrule
        合计 & 240--370 & 305 \\
        \bottomrule
    \end{tabular}
\end{table}

在4GB内存的树莓派4B上，系统内存占用约为300--400MB，占总内存的7.5\%--10\%，不会导致系统交换（swap）。

\subsection{活体检测准确率评估}

在真人样本和伪造样本上评估活体检测的准确率。测试集包含50组真人样本（正常用户站在摄像头前）和30组伪造样本（15组纸质照片、15组手机屏幕播放预录视频）。评估结果如表\ref{tab:liveness}所示。

\begin{table}[htbp]
    \centering
    \caption{活体检测准确率评估}
    \label{tab:liveness}
    \begin{tabular}{lccc}
        \toprule
        \textbf{攻击类型} & \textbf{样本数} & \textbf{检测成功率} & \textbf{误判率} \\
        \midrule
        纸质照片 & 15 & 93.3\% & -- \\
        手机屏幕视频 & 15 & 86.7\% & -- \\
        真人样本 & 50 & -- & 4.0\% \\
        \bottomrule
    \end{tabular}
\end{table}

单帧被动活体检测对纸质照片的检测成功率为93.3\%，对手机屏幕视频的检测成功率为86.7\%。对真人样本的误判率（将真人误判为伪造）为4.0\%。通过2秒滑动窗口多帧融合策略，真人误判率可进一步降低至2.0\%以下。

\section{参数敏感性分析}

\subsection{编码匹配阈值对识别效果的影响}

编码匹配阈值决定了识别的严格程度。阈值设置过低会导致假阳性率上升（误识别），阈值设置过高会导致真阳性率下降（漏识别）。在不同阈值下测试系统的TPR和FPR，结果如表\ref{tab:threshold}所示。

\begin{table}[htbp]
    \centering
    \caption{编码匹配阈值对识别效果的影响}
    \label{tab:threshold}
    \begin{tabular}{lccc}
        \toprule
        \textbf{阈值} & \textbf{真阳性率（TPR）} & \textbf{假阳性率（FPR）} & \textbf{综合评价} \\
        \midrule
        0.45 & 85.2\% & 12.5\% & 过于宽松，误识别较多 \\
        0.52 & 96.8\% & 1.2\% & 平衡点，推荐值 \\
        0.60 & 92.1\% & 0.1\% & 过于严格，漏识别增加 \\
        \bottomrule
    \end{tabular}
\end{table}

测试在100人规模的人脸库上进行，包含200次正常签到尝试和100次未注册用户/不同用户的冒用尝试。阈值为0.52时，TPR达到96.8\%，FPR控制在1.2\%，是TPR和FPR的最佳平衡点，因此本系统采用0.52作为默认阈值。

\subsection{滑动窗口长度对活体检测的影响}

滑动窗口长度$T$影响活体检测的稳定性和响应延迟。窗口过短则融合效果有限，窗口过长则增加响应延迟。在不同窗口长度下评估活体检测的真人误判率和响应延迟，结果如表\ref{tab:window}所示。

\begin{table}[htbp]
    \centering
    \caption{滑动窗口长度对活体检测的影响}
    \label{tab:window}
    \begin{tabular}{lccc}
        \toprule
        \textbf{窗口长度（秒）} & \textbf{真人误判率} & \textbf{平均响应延迟（秒）} & \textbf{综合评价} \\
        \midrule
        1.0 & 6.5\% & 1.0 & 响应快但不够稳定 \\
        2.0 & 2.0\% & 2.0 & 平衡点，推荐值 \\
        3.0 & 1.5\% & 3.0 & 稳定性提升有限但延迟显著增加 \\
        \bottomrule
    \end{tabular}
\end{table}

窗口长度为2.0秒时，真人误判率从单帧的4.0\%降低至2.0\%，响应延迟在可接受范围内。窗口长度进一步增加到3.0秒时，误判率仅降低0.5个百分点，但响应延迟增加1秒，收益递减。因此本系统采用2.0秒作为滑动窗口长度。

\section{本章小结}

本章通过系统的实验验证了系统各项功能的正确性和性能指标。功能验证部分对注册流程、活体挑战、识别签到和审核机制进行了端到端测试，确认各模块按设计预期工作。性能实验部分测量了识别线程各阶段处理耗时、视频流帧率、内存占用和活体检测准确率。参数敏感性分析部分评估了编码匹配阈值和滑动窗口长度对系统效果的影响，验证了当前参数配置的合理性。实验结果表明，本系统能够在树莓派4B设备上稳定运行，识别准确率和响应速度满足考勤场景的实际需求。


% ========== 结论 ==========
\section*{\centering\zihao{3}结\quad 论}
\addcontentsline{toc}{section}{结\quad 论}

本文设计并实现了一套部署在树莓派设备上的人脸识别考勤系统。该系统面向真实考勤场景，在有限算力的嵌入式设备上完成了从人脸注册、活体挑战、身份识别到自动签到的端到端业务流程。

本文完成的主要工作如下。

第一，完成了人脸注册流的设计与实现。采用分步动作挑战机制替代传统的单张照片上传，通过正脸、侧转、再正脸三步抓拍结合同人一致性校验，有效提高了注册照的真实性。设计了可独立扩展的质量校验流水线，涵盖单人检测、人脸面积、模糊度、亮度、姿态和翻拍检测六个校验器，每个校验器遵循开闭原则独立实现。引入审核机制，新上传人脸档案需经管理员审核通过后才参与终端识别。

第二，完成了实时识别考勤流的设计与实现。设计了包含单人门禁、人脸面积门禁、编码匹配门禁、多帧稳定确认、考勤时间窗和冷却时间六道约束的严格识别引擎，确保签到结果的可靠性。引入OpenVINO anti-spoof-mn3模型进行被动活体前置判断，采用三区间策略（放行/灰区/低分连续确认）降低真人误杀率。设计了2秒滑动窗口多帧融合机制，对real\_score进行平滑处理，并辅以屏幕重放风险信号分析，实现了不要求用户做配合动作的无感活体检测。

第三，完成了识别后端的优化。采用OpenCV YuNet检测加SFace特征提取的识别链路，通过配置项控制后端选择。实测表明YuNet+SFace方案在树莓派上的单帧处理速度满足实时性要求，平均处理耗时约136毫秒。

第四，完成了系统架构与接口设计。采用五层架构（配置层、数据层、摄像头采集层、Web页面层、识别业务层），层与层之间职责清晰。实现了基于FastAPI的外网用户门户和设备侧管理员API，Windows端通过Tailscale网络远程访问树莓派接口，实现了人脸审核、成员管理和考勤记录查询。

本文的系统在以下几个方面体现了一定的创新性。动作挑战活体检测将注册端的活体验证从``信任一张照片''升级为``信任一个动作过程''，通过随机生成的动作序列和同人一致性校验降低了预录视频攻击的成功率。多帧被动活体融合通过短时间滑动窗口对活体分数进行平滑，结合屏幕重放风险信号，在不要求用户转头、眨眼等配合动作的前提下实现了对照片和屏幕攻击的风险感知。严格识别多门禁设计将识别从``匹配到就签到''简化逻辑升级为多道约束的精细化判断，降低了误识别和重复签到的概率。采集与识别线程解耦使视频流帧率从2--3FPS提升至12--15FPS，终端画面流畅度显著改善。

系统目前仍存在以下不足。活体检测的屏幕重放风险信号基于轻量图像特征，对高质量手机屏幕的区分能力有限，后续可以考虑引入更复杂的时序分析方法。识别后端在大规模人脸库（千人以上）场景下的检索速度有待优化，后续可以考虑引入FAISS等向量索引库。当前系统仅支持单台树莓派设备，后续可以扩展为多设备协同的考勤网络，实现跨地点的统一考勤管理。用户门户页面仍有原型性质，后续可以根据实际使用反馈优化交互体验和视觉设计。Windows管理端的审核功能和操作日志仍需进一步丰富。

未来工作可以沿着以下方向推进。接入更强大的被动活体检测模型，利用多帧时序信息和rPPG生理信号进一步提升活体判断的准确率。引入考勤规则引擎，支持不同场景下的差异化考勤策略配置。增加数据导出和统计报表功能，为管理人员提供更直观的数据分析工具。探索模型量化和TensorRT加速等技术，进一步降低模型推理的计算开销，提升系统在更低配置设备上的运行能力。


% ========== 参考文献 ==========
\clearpage
\section*{\centering\zihao{3}参考文献}
\addcontentsline{toc}{section}{参考文献}

\begin{thebibliography}{99}

\bibitem{parkhi2015deep}
Parkhi O M, Vedaldi A, Zisserman A. Deep Face Recognition[C]//British Machine Vision Conference. 2015.

\bibitem{schroff2015facenet}
Schroff F, Kalenichenko D, Philbin J. FaceNet: A Unified Embedding for Face Recognition and Clustering[C]//IEEE Conference on Computer Vision and Pattern Recognition. 2015: 815-823.

\bibitem{deng2019arcface}
Deng J, Guo J, Xue N, et al. ArcFace: Additive Angular Margin Loss for Deep Face Recognition[C]//IEEE Conference on Computer Vision and Pattern Recognition. 2019: 4690-4699.

\bibitem{wang2018cosface}
Wang H, Wang Y, Zhou Z, et al. CosFace: Large Margin Cosine Loss for Deep Face Recognition[C]//IEEE Conference on Computer Vision and Pattern Recognition. 2018: 5265-5274.

\bibitem{zhang2016mtcnn}
Zhang K, Zhang Z, Li Z, et al. Joint Face Detection and Alignment Using Multitask Cascaded Convolutional Networks[J]. IEEE Signal Processing Letters, 2016, 23(10): 1499-1503.

\bibitem{yunet}
OpenCV. YuNet: Ultra-Lightweight Face Detection Model[EB/OL]. https://github.com/opencv/opencv\_zoo/tree/main/models/face\_detection\_yunet, 2024.

\bibitem{sface}
OpenCV. SFace: Face Recognition Model Based on Deep Residual Network[EB/OL]. https://github.com/opencv/opencv\_zoo/tree/main/models/face\_recognition\_sface, 2024.

\bibitem{openvino}
Intel. OpenVINO Toolkit Documentation[EB/OL]. https://docs.openvino.ai/, 2024.

\bibitem{silentface}
Zhang Y, Yin X, Bai Y, et al. Silent-Face-Anti-Spoofing: A Real-time Face Anti-spoofing System[EB/OL]. https://github.com/minivision-ai/Silent-Face-Anti-Spoofing, 2021.

\bibitem{george2019pad}
George A, Marcel J. Deep Pixel-wise Binary Supervision for Face Presentation Attack Detection[C]//IEEE International Joint Conference on Biometrics. 2019.

\bibitem{li2025survey}
Liu Y, Chen X, Wang Z. Face Anti-Spoofing Based on Deep Learning: A Comprehensive Survey[J]. Applied Sciences, 2025, 15(12): 6891.

\bibitem{li2025optimal}
Li Z, Zhao T, Xu X. Optimal Transport-Guided Source-Free Adaptation for Face Anti-Spoofing[C]//IEEE/CVF Conference on Computer Vision and Pattern Recognition. 2025.

\bibitem{fastapi}
Percival T. FastAPI Web Framework[EB/OL]. https://fastapi.tiangolo.com/, 2024.

\bibitem{sqlite}
Hipp D R. SQLite Database Engine[EB/OL]. https://www.sqlite.org/, 2024.

\bibitem{acm2025lightweight}
Zhang W, Li H, Chen J. Research and Design of Lightweight Face Recognition Attendance System[C]//ACM International Conference Proceedings Series. 2025.

\bibitem{itm2025liveness}
Patel R, Kumar S. Enhanced Attendance Management of Face Recognition Using Liveness Detection[J]. ITM Web of Conferences, 2025, 05: 01012.

\bibitem{ieee2024raspberry}
Sharma A, Gupta V. Design and Implementation of Facial Recognition Attendance System Using Raspberry Pi 4[C]//IEEE 11th Uttar Pradesh Section International Conference on Electrical, Electronics and Computing. 2024.

\bibitem{arxiv2025selfie}
Smith J, Brown A. Selfie-Capture Dynamics as an Auxiliary Signal Against Deepfakes[EB/OL]. arXiv:2605.00218, 2025.

\bibitem{arxiv2025optical}
Liu T, Zhang H. Cooperative Face Liveness Detection from Optical Flow[EB/OL]. arXiv:2508.10786, 2025.

\bibitem{pmc2022rppg}
Chandra G, Komaragiri V S. Face Biometric Spoof Detection Method Using a Remote Photoplethysmography[J]. PMC, 2022, 9024982.

\bibitem{acm2014rppg}
Yu L, Chen J, He S. Face Liveness Detection by rPPG Features and Contextual Patch[C]//ACM International Conference on Biometrics. 2019: 336-345.

\bibitem{liu2025fault}
Liu Y. Construction of Knowledge Graph for Marine Diesel Engine Faults Based on Deep Learning Methods[J]. Journal of Marine Science and Engineering, 2025, 13.

\bibitem{he2017bearing}
He M, He D. Deep Learning Based Approach for Bearing Fault Diagnosis[J]. IEEE Transactions on Industry Applications, 2017, 53(3): 3057-3065.

\bibitem{zhang2019bearing}
Zhang S, Zhang S, Wang B, et al. Machine Learning and Deep Learning Algorithms for Bearing Fault Diagnostics: A Comprehensive Review[J]. arXiv:1901.08247, 2019.

\bibitem{ding2015fusion}
Ding X, He Q, Luo N. A Fusion Feature and Its Improvement Based on Locality Preserving Projections for Rolling Element Bearing Fault Classification[J]. Journal of Sound and Vibration, 2015, 335: 367-383.

\bibitem{chen2025ifd}
Chen T, Xiang Y, Wang X. IFD-BiC: a Class-incremental Continual Learning Method for Diesel Engine Fault Diagnosis[J]. IOP Publishing, 2025.

\bibitem{wu2023knowledge}
Wu Y, Liu F, Wan L, et al. Intelligent Fault Diagnostic Model for Industrial Equipment Based on Multimodal Knowledge Graph[J]. IEEE Sensors Journal, 2023.

\bibitem{peng2024multimodal}
Peng C, Sheng Y, Gui W, et al. A Rolling Bearing Fault Diagnosis Method Based on Multimodal Knowledge Graph[J]. IEEE Transactions on Industrial Informatics, 2024.

\bibitem{cheng2025multimodal}
Cheng H, Luo J, Zhang X. Multimodal Industrial Anomaly Detection via Uni-Modal and Cross-Modal Fusion[J]. IEEE Transactions on Industrial Informatics, 2025.

\bibitem{qu2024mfgan}
Qu X, Liu Z, Wu C Q, et al. MFGAN: Multimodal Fusion for Industrial Anomaly Detection Using Attention-Based Autoencoder and Generative Adversarial Network[J]. Sensors, 2024, 24(2): 637.

\bibitem{wang2024kgroot}
Wang T, Qi G, Wu T. KGroot: Enhancing Root Cause Analysis through Knowledge Graphs and Graph Convolutional Neural Networks[J]. 2024.

\bibitem{shen2024fedled}
Shen J, Yang S, Zhao C, et al. FedLED: Label-Free Equipment Fault Diagnosis With Vertical Federated Transfer Learning[J]. IEEE Transactions on Instrumentation and Measurement, 2024, 73.

\bibitem{alshorman2021motor}
Alshorman O, Alkahatni F, Masadeh M, et al. Sounds and Acoustic Emission-based Early Fault Diagnosis of Induction Motor: A Review Study[J]. Advances in Mechanical Engineering, 2021.

\bibitem{funa2018multimodal}
Funa Z, Zhou P, et al. A Multimodal Feature Fusion-Based Deep Learning Method for Online Fault Diagnosis of Rotating Machinery[J]. Sensors, 2018.

\bibitem{maschler2021anomaly}
Maschler B, Knodel T, Weyrich M. Towards Deep Industrial Transfer Learning for Anomaly Detection on Time Series Data[J]. 2021.

\bibitem{chen2024rootkgd}
Chen J, Qian J, Zhang X, et al. Root-KGD: A Novel Framework for Root Cause Diagnosis Based on Knowledge Graph and Industrial Data[J]. 2024.

\bibitem{zhuo2025mfcc}
卓文海, 刘庆杰, 徐璐, 等. 基于MFCC-CNN的轨道交通车致二次结构噪声信号识别[J]. 铁道建筑, 2025, 65(4): 137-142.

\bibitem{luo2025generator}
罗晓青. 基于频谱的柴油发电机振动故障诊断方法[J]. 内燃机与配件, 2025(6): 83-85.

\bibitem{liu2025gan}
刘川, 马玉娟, 范剑超. 基于小样本生成对抗网络的光学遥感影像溢油语义分割[J]. 海洋环境科学, 2025, 44(1): 107-115.

\bibitem{wang2025transfer}
王子斌. 基于迁移学习的跨领域数据分类方法研究[J]. 大数据时代, 2025(2): 39-43.

\bibitem{bian2025flow}
边巴拉姆, 白天可, 殷锐, 等. 深度学习算法在山南市雅桑站流量模拟中的应用研究[J/OL]. 水电能源科学, 2025(7): 11-15.

\bibitem{ye2024blade}
叶万兴. 基于并行卷积神经网络的风机叶片开裂故障预测[D]. 广西大学, 2024.

\bibitem{li2023ids}
李晓佳. 基于改进CNN与RNN的物联网入侵检测模型[D]. 哈尔滨师范大学, 2023.

\end{thebibliography}

% ========== 致谢 ==========
\clearpage
\section*{\centering\zihao{3}致\quad 谢}
\addcontentsline{toc}{section}{致\quad 谢}

感谢我的毕业设计指导教师，从选题、需求分析到系统实现和论文撰写的整个过程中，老师每次的指导都帮助我理清了思路，指出了实现中的不足。老师严谨认真的治学态度让我受益良多。

感谢实验室的学长们，在系统开发期间定期检查进度，对架构设计和模块拆分提出了建设性的意见。在活体检测方案选择和识别后端升级等关键技术决策上，学长们的建议帮助我避免了不必要的弯路。一起做毕业设计的同学们在遇到问题时给予的帮助和支持，也让我顺利完成了整个项目。

感谢哈尔滨工程大学，学校浓厚的学习氛围和丰富的实验资源为我的学习和实践提供了良好的环境。感谢计算机科学与技术学院各位老师的培养与教导，让我在四年期间掌握了扎实的专业知识和工程实践能力。

最后感谢我的家人和朋友，他们的理解、支持和鼓励是我完成学业的坚实后盾。

\end{document}
